\section{Why the Predicted Distribution Collapses}
\label{sec:collapse}

\subsection{The observation}
\label{sec:collapse-observation}

We first reproduced the earlier problem in the new protocol. The probe on
$q_{\text{sem}}$ was trained on a random subset of 10k training scenes. Its
validation flow-matching loss is lowest at about 40k of 100k steps and then
rises, while the training error keeps falling. The validation dispersion drops
from 0.96 to 0.87, so the samples become narrower than the error justifies. On
the full 92k scenes the same drift exists but is much slower, and validation
still improves at 100k steps.

Before reading this as collapse, two other explanations had to be ruled out: the
sampler, and the size of the probe relative to the data.

\subsection{Part of the narrowing is the sampler}
\label{sec:euler}

Section~\ref{sec:flow-matching} showed that a single Euler step sends every
sample to the conditional mean. If 10 steps are still too coarse, part of this
effect remains, and the samples look narrower than the learned distribution.
Figure~\ref{fig:euler} confirms this on the best image-only probe (WM11-ctx). With
one step the dispersion is 0.37 and the mean error over samples looks best, while
best-of-32 is worst. The spread recovers with more steps, and the energy score
stops improving at 10 to 20 steps (2.88 at 10, 2.86 at 50).

\begin{figure}[htbp]
\centering
\includegraphics[width=\textwidth]{../pictures/euler_steps.png}
\caption{Effect of the number of Euler steps at sampling time, WM11-ctx probe,
9,675 validation scenes, 32 samples. Few steps pull the samples to the
conditional mean: the mean error over samples looks good, but best-of-32 and the
energy score are bad and the dispersion is far below 1.}
\label{fig:euler}
\end{figure}

We then re-evaluated every checkpoint of the 10k-scene run with 50 steps
(Figure~\ref{fig:collapse-curves}, blue). With 50 steps the dispersion goes only
from 1.03 to 0.95, and $\text{minADE}_{32}$ improves from 0.86 to 0.67\,m and
stays there. So for the small probe most of the narrowing we saw was a sampling
artefact, and the real narrowing during training is mild. This also means that a
model judged by mean ADE with few sampling steps will look unimodal whatever it
has learned.

\subsection{Real collapse needs capacity}
\label{sec:capacity}

The small probe has 19k parameters. A wider probe (hidden width 64, about 290k
parameters) on the same 10k scenes collapses for real, with any number of steps
(Figure~\ref{fig:collapse-curves}, red). At 6k steps it is the best
$q_{\text{sem}}$ model of all, with a validation energy of 5.65. Then
$\text{minADE}_{32}$ grows from 0.65 to 2.06\,m and the dispersion falls from 0.91
to 0.16, even with 50 steps. On training scenes $\text{minADE}_{16}$ goes to
0.14\,m. The model has stopped modelling a distribution and returns one
trajectory per scene, the future of a training scene that looks similar.
Figure~\ref{fig:collapse-samples} shows this. The fan of samples at 5k steps
becomes a single line at 100k, and in the first row that line turns left where
the car turned right.

The error of the sample mean does not show any of this. It stays the best of all
models (2.6\,m) during the whole collapse. A model selected by mean displacement
error would be the collapsed one, which is why we select on the energy score.

\begin{figure}[htbp]
\centering
\includegraphics[width=\textwidth]{../pictures/collapse_curves.png}
\caption{Validation metrics of the $q_{\text{sem}}$ probe on 10k training scenes
over training, re-evaluated on all 9,675 validation scenes with 32 samples. Blue:
small probe (width 16). Red: wide probe (width 64). Solid: 10 Euler steps,
dashed: 50. The wide probe collapses with either number of steps. The small
probe stays close to calibration with 50 steps.}
\label{fig:collapse-curves}
\end{figure}

\begin{figure}[htbp]
\centering
\includegraphics[width=0.78\textwidth]{../pictures/collapse_samples_d64.png}
\caption{Samples of the wide probe on validation scenes at 5k, 20k and 100k
training steps (blue: 32 samples, black: logged future, forward is up, units in
metres). The fan narrows to a single trajectory per scene, which is sometimes on
the wrong side (row 1).}
\label{fig:collapse-samples}
\end{figure}

Table~\ref{tab:capacity} shows that capacity moves the probe along one axis. More
capacity gives better early accuracy and an earlier and deeper collapse. A width-8
probe never leaves calibration, but it never becomes accurate either. On the full
92k scenes the wide probe still narrows (dispersion 0.67) but stays the most
accurate $q_{\text{sem}}$ probe. So reducing capacity prevents the collapse only
by underfitting. More data prevents it without that cost.

\begin{table}[htbp]
\centering
\caption{Probe capacity, $q_{\text{sem}}$ input, final values after 100k steps
on 2000 validation scenes with 16 samples and 10 Euler steps. Energy and minADE
in metres.}
\label{tab:capacity}
\begin{tblr}{colspec={llrrrrr}, row{1}={font=\bfseries}, cells={font=\small}}
\toprule
Width & Scenes & Params & Best energy (step) & Final energy & minADE & Dispersion \\
\midrule
8  & 10k & 5k   & 7.64 (100k) & 7.63 & 1.40 & 1.02 \\
16 & 10k & 19k  & 5.95 (48k)  & 6.04 & 0.91 & 0.87 \\
64 & 10k & 290k & 5.65 (6k)   & 8.33 & 2.46 & 0.05 \\
16 & 92k & 19k  & 5.56 (100k) & 5.55 & 0.79 & 0.93 \\
64 & 92k & 290k & 4.54 (42k)  & 4.69 & 0.78 & 0.67 \\
\bottomrule
\end{tblr}
\end{table}

\subsection{The onset of memorization grows with the data}
\label{sec:scaling}

If the collapse is memorization, theory makes a sharp prediction: the training
time at which it starts should grow linearly with the number of training scenes
$n$~\cite{bonnaire,favero}. We trained the same small probe on 2.5k, 5k, 10k,
20k, 46k and 92k scenes (Figure~\ref{fig:scaling}, Table~\ref{tab:scaling}).

The validation loss reaches its minimum at 8k, 12k and 40k steps for 2.5k, 5k and
10k scenes. For 20k scenes and more it does not turn within 100k steps. A
log-log fit over the three sizes where the minimum is reached gives a slope of
1.16 for the loss minimum, 1.13 for the best energy and 0.90 for the step where
the gap between training and validation energy starts to grow. All three are
close to the predicted slope of 1. The smaller the training set, the earlier the
validation loss turns and the further the dispersion falls below 1, down to 0.63
for 2.5k scenes. The extrapolated onset for 20k scenes is about 80k steps, just
beyond what the run shows.

So the law holds for a small conditional trajectory model on real driving data,
not only for the image and language models it was shown for. With only three
uncensored points the
slope is a rough number. But the order is the same in every metric: less data
means earlier and stronger collapse.

\begin{figure}[htbp]
\centering
\includegraphics[width=\textwidth]{../pictures/data_scaling.png}
\caption{Small $q_{\text{sem}}$ probe trained on $n$ scenes. From left: energy
score (solid validation, dotted training), growth of the validation-training
gap, validation dispersion, and the onset step against $n$ on log-log axes.
Crosses mark sizes where the onset is not reached within 100k steps.}
\label{fig:scaling}
\end{figure}

\begin{table}[htbp]
\centering
\caption{Onset of memorization against training-set size, small
$q_{\text{sem}}$ probe, 100k steps.}
\label{tab:scaling}
\begin{tblr}{colspec={rrrrr}, row{1}={font=\bfseries}}
\toprule
Scenes & Val loss minimum at & Final val energy & Final dispersion & Val / train minADE \\
\midrule
2,500  & 8k          & 8.55 & 0.63 & 1.57 / 0.65 \\
5,000  & 12k         & 6.59 & 0.82 & 1.00 / 0.68 \\
10,000 & 40k         & 6.04 & 0.87 & 0.91 / 0.75 \\
20,000 & not reached & 5.79 & 0.91 & 0.85 / 0.73 \\
46,000 & not reached & 5.66 & 0.92 & 0.83 / 0.75 \\
92,381 & not reached & 5.55 & 0.93 & 0.79 / 0.74 \\
\bottomrule
\end{tblr}
\end{table}

\subsection{A polynomial code does not prevent collapse}
\label{sec:codec}

A natural idea is to restrict the output to smooth trajectories, so that the
model cannot fit the training futures in detail. We compared polynomial codes of
order 2, 3 and 6 with the raw 24 numbers, on 10k scenes where the collapse is
visible (Table~\ref{tab:codec}).

\begin{table}[htbp]
\centering
\caption{Trajectory code, $q_{\text{sem}}$, small probe, 10k scenes, final
validation values.}
\label{tab:codec}
\begin{tblr}{colspec={lrrrr}, row{1}={font=\bfseries}}
\toprule
Code & Size & Val energy & Val minADE & Dispersion \\
\midrule
polynomial, order 2 & 6  & 6.00 & 0.93 & 0.82 \\
polynomial, order 3 & 9  & 6.04 & 0.91 & 0.87 \\
polynomial, order 6 & 18 & 8.25 & 1.58 & 1.10 \\
raw poses           & 24 & 7.07 & 2.17 & 1.14 \\
\bottomrule
\end{tblr}
\end{table}

Low-order polynomials are much better targets. Raw poses have no notion of
smoothness, so independent errors per pose give jagged paths, and the order-6
basis is badly conditioned, so small errors in the code become large errors in
the path. But the low-order codes still narrow below calibration, order 2 even
slightly more than order 3. Raw poses and order 6 avoid the collapse only because
they are not fitted well enough to collapse. A simpler target makes the
distribution easier to learn, and that includes the memorized one.

\subsection{Summary}

The narrowing of the predicted distribution has two causes. One is a measurement
artefact: coarse ODE integration pulls the samples to the mean. The other is
memorization, which needs a probe that can fit its training set and appears
later the more data there is. Neither the features nor the trajectory code cause
it. In practice it is avoided by training on all 92k scenes, sampling with at
least 10 to 20 steps, and selecting checkpoints on the validation energy score.
All results in the next sections follow these rules. Every probe there has a
dispersion between 0.88 and 0.95 on the validation logs. Lower values appear only
on inputs unlike the training data (Section~\ref{sec:generalization}).
